\documentclass[11pt]{article}
\usepackage{amsmath, amssymb, amsthm}
\usepackage{geometry}
\usepackage{graphicx}
\usepackage{subcaption}
\usepackage{setspace}
\usepackage{mathtools}
\usepackage{bm}
\usepackage{float}

\geometry{margin=1in}
\onehalfspacing

\title{When Do Correlation and Causal Forests Differ?}
\author{Martin Huber\\
\small {University of Fribourg, Department of Economics} \smallskip }

\begin{document}
\section{The underlying idea: multivariate fixed-set iterations}\label{sec:idea}

\subsection{The idea in the univariate case}\label{sec:idea_uni}

The intuition for identification of the univariate version of model  follows the idea from . The univariate model is the one where all variables, observable and unobseervable, are univariate and where both $m$ and $h$ satisfy the rank-invariance condition in their second argument, i.e.\ they are strictly increasing and continuous functions in $\varepsilon$ and $U$, respectively. 

$m(X,\varepsilon)$ in model  maps the distribution $F_\varepsilon$ of the unobservable to the conditional \textit{counterfactual} distribution $F_{Y(X)}$ for \textit{exogenous} $X$, i.e.\ for an $X$ with $X\perp\varepsilon$. The notation $F_{Y(X)}$ is Rubin's counterfactual notation : due to the endogeneity problem, i.e.\ the fact that $X$ and $\varepsilon$ are not independent, $F_{Y(X)}$ is unobservable and the observable distribution $F_{Y|X}$ does not coincide with $F_{Y(X)}$. In order to identify $m$, we would need to know $F_{Y(X)}$, i.e.\ we want to know how the model would behave for a change in $X$ that does not affect the distribution of $\varepsilon$.

The idea is to vary $Z$ in such a way that $X$ varies but $\varepsilon$ stays constant. The fact that this is possible if $Z$ is only binary was first shown in  and . The argument rests on the fact that using the binary instrument $Z$ with realizations $z$ and $z'$, there are two maps that do not change the distribution $F_\varepsilon$ of $\varepsilon$, but change values of $X$. This hence captures the exogenous effect of $X$ on $Y$. These two maps are depicted in Figure .

\begin{figure}[H]
    \centering
    \includegraphics[width=0.7\textwidth]{figures/fig1.pdf}
    \caption{Fixed-point iteration in a univariate framework.}
    \label{fig:fig1}
\end{figure}

The first map changes $z\mapsto z'$ for a fixed $x$, i.e.\ switches the distributions $F_{X|Z=z}(x)$ and $F_{X|Z=z'}(x)$ (the ``vertical'' map in Figure ). The fact that $F_{\varepsilon|X,Z}$ is not affected by this follows from the classical control variable approach  and because $Z$ has no effect on $m$ due to the exclusion restriction of $Z$. A control variable $V$ is such that $X\perp\varepsilon|V$ and can be constructed via $V=F_{X|Z}(X)$ . In particular, it contains the same information as the unobservable $U$. Hence, by conditioning on $V$ and the fact that $Z$ is independent of $(\varepsilon,U)$, the vertical shift changes $Z$ but does not change $X$, which therefore does not affect the function $m(x,\varepsilon)$ we want to identify .
\end{document}
